\section*{Overview}

Min-cost max-flow extends ordinary max flow by attaching a cost to each unit of flow on each edge. The goal is not
just to push flow, but to push it as cheaply as possible. This is the right tool when a plain feasibility or cardinality
answer is not enough.

\section*{When to Use It}

Use min-cost max-flow when:

\begin{itemize}
  \item the network model is already natural,
  \item each assignment or transfer has a linear cost,
  \item you need either the cheapest way to send a fixed amount or the cheapest maximum flow.
\end{itemize}

Contest examples include weighted assignment, transportation, balancing supply and demand, and extracting \(k\) cheapest
disjoint paths in small-to-medium graphs.

\section*{Core Idea}

The standard contest implementation is successive shortest augmenting path:

\begin{itemize}
  \item keep a residual network with capacities and edge costs,
  \item repeatedly find the cheapest augmenting path from \(s\) to \(t\),
  \item push as much flow as possible on that path,
  \item update residual capacities and accumulated cost.
\end{itemize}

Potentials reweight edges so Dijkstra can still be used even when the residual network contains negative-cost reverse
edges.

\section*{Key Insight}

The reverse edge is what makes cost corrections possible. If later decisions should undo some earlier expensive choice,
the residual network exposes that option as a negative-cost reverse edge.

\section*{Operations / Main Technique}

\begin{itemize}
  \item add forward and reverse residual edges,
  \item shortest path in the residual network with reduced costs,
  \item augment,
  \item update vertex potentials.
\end{itemize}

\paragraph{Worked example.}
In weighted bipartite assignment, each left-right compatibility edge carries a cost. Sending one unit of flow from
source to left vertex to right vertex to sink corresponds exactly to choosing one assignment pair.

\section*{Correctness Intuition}

Each augmentation chooses the cheapest possible additional residual path under the current reweighting. Potentials do
not change which paths are optimal; they only keep reduced edge costs nonnegative so Dijkstra remains valid. When no
more augmenting path exists, the flow is maximal, and the accumulated augmentations are minimum-cost among flows of that
value.

\section*{Complexity Analysis}

With Dijkstra and potentials, one augmentation is roughly \(O(E \log V)\). The total cost depends on how many times the
algorithm augments, so the method is best when the graph and required flow are not enormous.

\section*{Implementation}

The sample returns both:

\begin{itemize}
  \item the achieved flow,
  \item the total minimum cost.
\end{itemize}

That makes it usable both for ``send exactly \(k\)'' and ``find the cheapest maximum flow'' style tasks.

\section*{Common Pitfalls}

\begin{itemize}
  \item Forgetting to add the reverse edge with negated cost.
  \item Not updating potentials after each shortest-path phase.
  \item Using the algorithm on constraints where the number of augmentations is too large.
  \item Confusing minimum-cost maximum flow with minimum-cost feasible flow of a specified amount.
\end{itemize}

\section*{Variants / Extensions}

\begin{itemize}
  \item Min-cost flow for a fixed required amount.
  \item Lower bounds and circulation reductions.
  \item Potentials initialized by Bellman-Ford or SPFA if true negative edges exist initially.
\end{itemize}

\section*{Practice Problems}

\begin{itemize}
  \item Weighted bipartite matching through a flow model.
  \item Send \(k\) units through a graph with edge costs.
  \item Transportation between supply and demand nodes.
\end{itemize}

\section*{References}

\begin{itemize}
  \item \href{https://cp-algorithms.com/graph/min_cost_flow.html}{cp-algorithms: Min-Cost Flow}.
  \item \href{https://usaco.guide/adv/min-cost-flow}{USACO Guide: Min-Cost Flow}.
  \item \href{https://codeforces.com/catalog}{Codeforces Catalog}.
\end{itemize}
